% Replace the sample content below with the record of your own session.
\section{\lblGlance}
\begin{table}[H]
  \centering
  \begin{tabularx}{\linewidth}{@{}l@{\hspace{1.2cm}}X@{}}
    \toprule
    \metalabel{\lblDate} & \RepDate \\
    \metalabel{\lblDuration} & \RepDuration \\
    \metalabel{\lblAttendees} & Two members of the support team and the operations lead \\
    \metalabel{\lblFormat} & Video call with a shared screen \\
    \metalabel{\lblTopics} & AI agents in customer support, review of replies, handling of customer data \\
    \bottomrule
  \end{tabularx}
\end{table}

\begin{itemize}
  \item Agents can draft replies for the support team now, but should not send them alone.
  \item Customer data needs a clear boundary before any agent reads a ticket.
  \item A small pilot with measured results is a better start than a full rollout.
\end{itemize}

\section{\lblQuestions}
\begin{enumerate}
  \item Which support tasks can an AI agent do safely today?
  \item How do we stop an agent from sharing customer data it should not see?
  \item Who checks the agent's work, and how much time does that take?
  \item How will we know the agent is helping and not just fast?
\end{enumerate}

\section{\lblFound}
\subsection{Tickets mix routine and sensitive requests}
\findingpart{\lblContext}
The team handles about 300 tickets a week. Most are order status and password
questions; a few involve refunds, complaints or personal data.

\findingpart{\lblObservation}
Routine tickets are a good fit for an agent. Sensitive ones are mixed into the
same queue with no label, so an agent cannot tell them apart.

\subsection{Replies are not reviewed in a consistent way}
\findingpart{\lblContext}
Today each team member checks replies in their own way, and there is no written
standard for tone or accuracy.

\findingpart{\lblObservation}
Without a shared standard, a reviewer cannot say whether a draft is good, and
the team cannot measure the agent against anything.

\begin{insight}
  The safe limit for an agent is set by what it can see and what it can send,
  not by how clever it is. Fix those two limits first.
\end{insight}

\begin{mermaid}[style=editorial,width=0.9\linewidth,
    caption={Proposed handling of a ticket},pos=H]
  flowchart LR
  A[Ticket] --> B[Agent drafts] --> C[Person reviews] --> D[Send]
\end{mermaid}

\section{\lblRecommendations}
\begin{table}[H]
  \centering
  \begin{tabularx}{\linewidth}{@{}lXll@{}}
    \toprule
    & \sffamily\small\bfseries\lblRecommendation & \sffamily\small\bfseries\lblPriority &
    \sffamily\small\bfseries\lblEffort \\
    \midrule
    \recrow{1}{Keep a person between the agent and the customer}{Drafts only; a person sends every reply.}{high}{small}
    \recrow{2}{Hide personal data from the agent}{Mask names, emails and payment details before the ticket is read.}{high}{medium}
    \recrow{3}{Write a one-page reply standard}{Gives reviewers and the pilot a shared measure of quality.}{medium}{small}
    \recrow{4}{Run a four-week pilot on routine tickets}{Measure time saved and errors found before widening the scope.}{medium}{medium}
    \recrow{5}{Add a monthly review of agent mistakes}{Keeps the rules current as the tickets change.}{low}{small}
    \bottomrule
  \end{tabularx}
\end{table}

\section{\lblResources}
\begin{table}[H]
  \centering
  \begin{tabularx}{\linewidth}{@{}Xl@{}}
    \toprule
    \sffamily\small\bfseries\lblResource & \sffamily\small\bfseries\lblLink \\
    \midrule
    \resrow{Checklist for reviewing an agent's replies}{https://example.org/checklist}
    \resrow{Template for a one-page reply standard}{https://example.org/reply-standard}
    \resrow{Guide to masking personal data in tickets}{https://example.org/masking}
    \resrow{Pilot plan worksheet}{https://example.org/pilot-plan}
    \bottomrule
  \end{tabularx}
\end{table}

\section{\lblNextSteps}
\begin{table}[H]
  \centering
  \begin{tabularx}{\linewidth}{@{}Xlr@{}}
    \toprule
    \sffamily\small\bfseries\lblAction & \sffamily\small\bfseries\lblOwner &
    \sffamily\small\bfseries\lblTarget \\
    \midrule
    \actrow{Label sensitive tickets in the queue}{client}{20 November 2026}
    \actrow{Send the masking guide and an example setup}{consultant}{23 November 2026}
    \actrow{Draft the one-page reply standard}{client}{27 November 2026}
    \actrow{Review the draft standard and the pilot plan}{consultant}{2 December 2026}
    \actrow{Start the four-week pilot}{client}{7 December 2026}
    \bottomrule
  \end{tabularx}
\end{table}

\section{\lblNotes}
\begin{rawnotes}
  \begin{itemize}
    \item Team worried about wrong refund promises; agent must never state a refund.
    \item Current tool keeps full ticket history, including card digits in free text.
    \item Operations lead wants a weekly number: tickets answered, edits per draft.
    \item Idea: start with order status only, add password help later.
  \end{itemize}
  Rule sketch agreed in the call:
  \begin{code}[lang=yaml]
agent:
  can_read: [order_status, shipping]
  can_send: false
  blocked: [refund, complaint, payment]
  \end{code}
\end{rawnotes}

\section{\lblFollowUp}
{\raggedright\hyphenpenalty=10000\exhyphenpenalty=10000 \RepBooking\par}
